% Festlegung des allgemeinen Dokumentenformats
\documentclass[a4paper,12pt,headsepline]{article}

% Schrift
\usepackage[T1]{fontenc}
\usepackage{lmodern}
\usepackage[utf8]{inputenc}
\usepackage[ngerman]{babel}

% Bilder
\usepackage{graphicx}

% Variablen
% Variablen
\newcommand{\titleDocument}{Bachelorarbeit}
\newcommand{\subjectDocument}{im Studiengang Informatik}
\newcommand*{\bildquelle}{
  \footnotesize Quelle:
}
\newcommand{\code}[1]{\noindent\ignorespaces\texttt{#1}}

% Code
\usepackage{listings}
\usepackage{color}

\definecolor{black}{rgb}{0,0,0}
\definecolor{dkgreen}{rgb}{0,0.6,0}
\definecolor{gray}{rgb}{0.5,0.5,0.5}
\definecolor{mauve}{rgb}{0.58,0,0.82}

\lstset{literate=%
    {Ö}{{\"O}}1
    {Ä}{{\"A}}1
    {Ü}{{\"U}}1
    {ß}{{\ss}}1
    {ü}{{\"u}}1
    {ä}{{\"a}}1
    {ö}{{\"o}}1
    {~}{{\textasciitilde}}1
}

\lstdefinestyle{Bash} {
  frame=single,
  language=Bash,
  aboveskip=3mm,
  belowskip=3mm,
  showstringspaces=false,
  columns=flexible,
  basicstyle={\small\ttfamily},
  numbers=none,
  numberstyle=\tiny\color{black},
  keywordstyle=\color{black},
  commentstyle=\color{black},
  stringstyle=\color{black},
  breaklines=true,
  breakatwhitespace=true,
  tabsize=3
}

\lstset{emph={\$},emphstyle=\textbf}

\lstdefinestyle{Python} {
  frame=single,
  language=Bash,
  aboveskip=3mm,
  belowskip=3mm,
  showstringspaces=false,
  columns=flexible,
  basicstyle={\small\ttfamily},
  numbers=left,
  numberstyle=\tiny\color{black},
  keywordstyle=\color{black},
  commentstyle=\color{gray},
  stringstyle=\color{gray},
  breaklines=true,
  breakatwhitespace=true,
  tabsize=2
}

% mehrseitige Tabellen ermöglichen
\usepackage{longtable}

% Packet für Seitenrandabstände und Einstellung für Seitenränder
\usepackage{geometry}
\geometry{left=3.5cm, right=2.5cm, top=2.5cm, bottom=2cm}

% bricht lange URLs "schön" um
\usepackage[hyphens,obeyspaces,spaces]{url}

% Festlegung Art der Zitierung - Havardmethode: Abkuerzung Autor + Jahr
\bibliographystyle{alphadin}

% Paket für Zeilenabstand
\usepackage{setspace}
\onehalfspacing

% für Bildbezeichner
\usepackage{capt-of}

% für Stichwortverzeichnis
\usepackage{makeidx}

% Für Phantomsection
\usepackage{hyperref}

% Konfiguriere das Inhaltsverzeichnis
\usepackage{tocbasic}
\DeclareTOCStyleEntries[
  raggedentrytext,
  numwidth=0pt,
  numsep=1ex,
  dynnumwidth,
]{tocline}{chapter,section,subsection,subsubsection,paragraph,subparagraph}
\DeclareTOCStyleEntries[
  indent=0pt,
  linefill=\TOCLineLeaderFill,
]{tocline}{section,subsection,subsubsection,paragraph,subparagraph}

% Titel
\title{Bachelorarbeit}

% Autor
\author{Andreas Huber}

% Datum
\date{\today}

%
% Start
% des
% Dokuments
%
\begin{document}

% Titelseite
\thispagestyle{empty}

\begin{figure}[t]
 \centering
 \includegraphics[width=0.4\textwidth]{assets/oth/logo}
\end{figure}

\begin{verbatim}
\end{verbatim}

\begin{center}
    \Large{Ostbayerische Technische Hochschule Regensburg} \\
    \Large{Fakultät für Informatik und Mathematik}
\end{center}

\begin{verbatim}
\end{verbatim}

\begin{center}
    \doublespacing
    \textbf{\huge{\titleDocument}}\\

    \onehalfspacing

    \begin{center}
        Zur Erlangung des akademischen Grades \\ Bachelor of Science (B. Sc.)
    \end{center}

    \begin{verbatim}
    \end{verbatim}

    \begin{doublespace}
        \textbf{\Large{{~\subjectDocument}}}
    \end{doublespace}
\end{center}

\begin{verbatim}
\end{verbatim}

\begin{verbatim}
\end{verbatim}

\begin{flushleft}
    \begin{tabularx}{\linewidth}{@{}>{\bfseries}l@{\hspace{.9em}}X@{}}
        \textbf{Vorgelegt von:} & Andreas Huber <andreas.huber@st.oth-regensburg.de> \\
        \textbf{Matrikelnummer:}& 3180161 \\
        \textbf{Studiengang:}   & Allgemeine Informatik \\
                                & \\
        \textbf{Erstgutachter:} & Prof. Dr. Markus Heckner \\
        \textbf{Zweitgutachter:}& Prof. Dr. Johannes Schildgen \\
                                & \\
        \textbf{Abgabefrist:}   & 25. April 2022 \\
    \end{tabularx}
\end{flushleft}
\newpage

% Römische Nummerierung
\pagenumbering{roman}

% Inhaltsverzeichnis
\tableofcontents
\newpage

% Arabische Seitennummerierung ab hier
\pagenumbering{arabic}

% Einleitung
\section{Einleitung}\label{einleitung}

Lorem ipsum dolor sit amet, consetetur sadipscing elitr, sed diam nonumy eirmod tempor invidunt ut labore et dolore magna aliquyam erat, sed diam voluptua. At vero eos et accusam et justo duo dolores et ea rebum. Stet clita kasd gubergren, no sea takimata sanctus est Lorem ipsum dolor sit amet. Lorem ipsum dolor sit amet, consetetur sadipscing elitr, sed diam nonumy eirmod tempor invidunt ut labore et dolore magna aliquyam erat, sed diam voluptua. At vero eos et accusam et justo duo dolores et ea rebum. Stet clita kasd gubergren, no sea takimata sanctus est Lorem ipsum dolor sit amet.
\newpage

% Römische Nummerierung
\pagenumbering{roman}
\setcounter{page}{3}

% Literaturliste soll im Inhaltsverzeichnis auftauchen
\phantomsection
\addcontentsline{toc}{section}{Quellenverzeichnis}
% Literaturverzeichnis anzeigen
\renewcommand\refname{Quellenverzeichnis}
\printbibliography
% \bibliography{assets/literatur}
\newpage

\end{document}